% Names of the subsection reflect the content that is required. More subsections can be added if required
\chapter{Literature Survey}
\label{section:lit_survey}

Object tracking in video is commonly formulated as a detection-and-association pipeline. Detection 
localizes objects in each frame, while tracking predicts motion and preserves temporal consistency. 
The overall quality depends on both stages and on implementation efficiency, especially for real-time 
or hardware-constrained deployments \cite{Milan2016MOT16, Bewley2016_sort}. This chapter reviews prior 
work that is most relevant to the present pipeline, where foreground segmentation is used for detection 
and a SORT-style tracker is used for online state estimation.

\section{Related Work}
\label{section:related_work}

\subsection{Detection Methods for Video Tracking}
Early detection approaches relied on handcrafted features and classical classifiers. Viola--Jones used 
Haar-like features with cascade classification for fast detection \cite{ViolaJones2001}, and Dalal--Triggs 
showed robust pedestrian detection using HOG descriptors \cite{DalalTriggs2005}. Deep learning methods 
later improved accuracy, including Faster R-CNN \cite{Ren2015FasterRCNN} and one-stage detectors such as 
YOLO \cite{Redmon2016YOLO}. While these methods are highly effective, they often require higher compute 
and memory budgets than lightweight motion-based pipelines.

For tracking tasks where the foreground object is visually prominent and scene dynamics are manageable, 
background subtraction remains an efficient alternative to heavy per-frame object detectors. This motivates 
a foreground-driven detection stage for the current project.

\subsection{Background Subtraction and Adaptive Gaussian Models}
Background subtraction is a widely used strategy for extracting moving objects in surveillance and 
scene monitoring. Stauffer and Grimson introduced adaptive Gaussian mixture modeling for per-pixel 
background representation under dynamic scene conditions \cite{StaufferGrimson1999}. Zivkovic proposed 
adaptive updates that improve model efficiency and robustness in practical 
settings \cite{Zivkovic2004, ZivkovicHeijden2006}. These foundations are reflected in common computer 
vision toolkits such as OpenCV, where MOG2 is used for real-time foreground mask 
estimation \cite{Bradski2000OpenCV}.

In practice, background subtraction is typically followed by denoising and connected component 
or contour analysis to convert binary masks into bounding boxes. This combination offers low latency 
and deterministic behavior, which are useful when interfacing with hardware-friendly tracking stages.

\subsection{Tracking-by-Detection and Data Association}
Tracking-by-detection decomposes MOT into prediction, association, and correction. The Kalman 
filter remains a standard choice for linear state estimation due to recursive updates and low computational 
cost \cite{Kalman1960}. In general MOT settings, the Hungarian algorithm is commonly used for optimal 
bipartite assignment between predicted tracks and current detections \cite{Kuhn1955Hungarian}.

Community benchmarks such as MOTChallenge have standardized tracker evaluation and shown that detector 
quality and association strategy strongly influence final tracking metrics \cite{Milan2016MOT16}. These 
findings justify evaluating both geometric overlap and temporal consistency when comparing tracker variants.

\subsection{SORT and Successor Approaches}
SORT (Simple Online and Realtime Tracking) combines Kalman prediction with IoU-based association and 
Hungarian matching \cite{Bewley2016_sort}. Its main advantage is a strong speed-accuracy balance for 
online use, with minimal appearance modeling overhead. DeepSORT extends SORT by adding appearance embeddings 
to reduce identity switches \cite{Wojke2017DeepSORT}, and ByteTrack improves association robustness by using 
low-confidence detections in a second matching stage \cite{Zhang2022ByteTrack}.

Although these extensions improve identity handling in crowded scenes, they increase algorithmic and 
implementation complexity. For a single-target or dominant-target pipeline, association can be 
simplified using IoU gating without full multi-track optimization.

\subsection{Hardware-Aware Tracking Implementations}
Edge and embedded deployments require attention to latency, throughput, and resource usage in addition 
to accuracy. Prior surveys emphasize algorithm-hardware co-design for efficient real-time 
systems \cite{Sze2017EfficientDNN}. Tracking kernels based on matrix and vector arithmetic are 
especially suitable for hardware mapping. Recent work on FPGA-oriented Kalman tracking demonstrates 
practical throughput gains and validates this direction for motion tracking systems \cite{BABU2022101084}.

\subsection{Research Gap and Project Positioning}
The reviewed literature suggests three practical gaps for deployment-focused tracking systems:
\begin{enumerate}
	\item Many high-accuracy pipelines rely on compute-intensive detection or appearance modeling 
	that is expensive for constrained implementations.
	\item Background-subtraction-based pipelines are efficient, but are less frequently integrated 
	with explicit software-to-hardware tracker validation workflows.
	\item Single-object tracking use cases can simplify association logic, but this simplification 
	is often not described in a hardware implementation context.
\end{enumerate}

Accordingly, this project adopts a foreground-segmentation detector with SORT-inspired state estimation, 
then compares software reference outputs with custom hardware-oriented tracking outputs. The objective 
is to maintain online behavior while reducing implementation complexity.

\section{Tools and Technologies}
\label{section:tools_tech}

\subsection{Foreground Detection Stack}
The detection front-end is based on adaptive background subtraction (MOG2) followed by contour-based 
localization. This choice aligns with literature on low-latency motion segmentation in video scenes 
\cite{StaufferGrimson1999, Zivkovic2004, ZivkovicHeijden2006}. OpenCV provides optimized implementations 
for mask generation, filtering, contour extraction, and visualization \cite{Bradski2000OpenCV}.

\subsection{Tracking Core}
The tracking core follows SORT principles: linear state prediction with a Kalman filter and 
geometric validation through IoU \cite{Kalman1960, Bewley2016_sort}. Canonical SORT uses Hungarian 
assignment in multi-object settings \cite{Kuhn1955Hungarian}; however, the current implementation 
focuses on single-object tracking and therefore uses simplified association logic based on gating.

\subsection{Software Prototyping and Reference Comparison}
Python-based prototyping is used to generate detections, run a reference SORT flow, and produce 
frame-aligned outputs for analysis. This supports reproducible comparison of custom tracker behavior 
against a known software baseline.

\subsection{Hardware Mapping Perspective}
The design focus is to map deterministic arithmetic-heavy tracking blocks to hardware-friendly modules 
while preserving online updates. This perspective is consistent with prior real-time and FPGA-oriented 
tracking literature \cite{Sze2017EfficientDNN, BABU2022101084}.

\subsection{Summary of Survey}
The survey supports a practical architecture for this project: efficient foreground-based detection, 
lightweight online tracking, and systematic software-hardware comparison. This framing informs the 
methodology chapter and motivates the chosen trade-off between tracking robustness and implementation 
simplicity.